% answers_short.tex —— 2011 年·北京卷·文科 · 简短答案速查
% 由 tools/make_short_answers.py 自 answers.tex 的 \gksolhead/\gkanswer 提取，
% 并与 2011/2011_数学真题_北京文_答案要点.md 逐题比对后生成，勿手改（改源 MD 后重跑）。

\begin{gknotice}
\textbf{使用说明}：本节只给最终答案便于快速核对，完整推导见解析版。答案为编写者独立推导，\textbf{非官方标准答案}。
\end{gknotice}

\gkpart{选择题}
\begin{center}\begin{tabular}{ccccccccc}
\toprule
题号 & 1 & 2 & 3 & 4 & 5 & 6 & 7 & 8 \\
\midrule
答案 & D & A & D & D & B & C & B & A \\
\bottomrule
\end{tabular}\end{center}

\gkpart{填空题}
\begin{itemize}[leftmargin=2.6em]
  \item[\textbf{9.}] $a=\dfrac{5\sqrt2}{3}$。
  \item[\textbf{10.}] $2$。
  \item[\textbf{11.}] $1$。
  \item[\textbf{12.}] $q=-2$；$a_1+a_2+\cdots+a_n=\dfrac{1-(-2)^n}{6}$。
  \item[\textbf{13.}] $(0,1)$。
  \item[\textbf{14.}] $N(0)=6$；$N(t)$ 的所有可能取值为 $6,7,8$。
\end{itemize}

\gkpart{解答题}
\begin{itemize}[leftmargin=2.6em]
  \item[\textbf{15.}] (1) 最小正周期 $T=\pi$；(2) 最大值 $2$（$x=\dfrac{\pi}{6}$），最小值 $-1$（$x=-\dfrac{\pi}{6}$）
  \item[\textbf{16.}] (1) 平均数 $\dfrac{35}{4}$，方差 $\dfrac{11}{16}$；(2) $\dfrac14$
  \item[\textbf{17.}] (1) 证明见解析；(2) 证明见解析；(3) 存在；该点是矩形 $DEFG$ 的中心（$DF$ 与 $EG$ 的交点，也是 $PC$、$AB$ 中点连线的中点）
  \item[\textbf{18.}] (1) 在 $(-\infty,k-1)$ 上单调递减，在 $(k-1,+\infty)$ 上单调递增；(2) $k\leqslant1$ 时最小值为 $-k$；$1<k<2$ 时最小值为 $-\mathrm{e}^{k-1}$；$k\geqslant2$ 时最小值为 $(1-k)\mathrm{e}$
  \item[\textbf{19.}] (1) $\dfrac{x^2}{12}+\dfrac{y^2}{4}=1$；(2) $\dfrac92$
  \item[\textbf{20.}] (1) $A_5:0,1,0,1,0$（相邻两项之差均满足要求，且 $a_1=a_3=0$）；(2) 证明见解析；(3) $n$ 的最小值为 $9$，如 $A_9:4,3,2,1,0,-1,-2,-3,-4$
\end{itemize}
